% Complete new low-endpoint calculation. Historical sources are unchanged.
\section{The original scalar Gamma endpoint to fifth-order finite error}
\label{sec:LER}

Retain the original density and both literal moment indices:
\[
 \sigma(y)=\frac{|\Gamma(1/4+iy/2)|^2}{2\pi},\qquad
 I_p=\int_0^\infty y^p\sigma(y)\,dy\quad(p\ge0),\qquad
 \mu_a=\int_{\mathbb R}y^{2a}\sigma(y)\,dy=2I_{2a}.
 \tag{LER1}
\]
Thus $\mu_a=m_{2a}$ in LET's all-power notation, and
$\mu_0=\sqrt{2\pi}$. The factors two and the original mass are kept
until cancellation in the specified ratios. Let $\alpha_\Gamma=\pi/2$ and
$L(y)=-\sigma'(y)/\sigma(y)$. The complete accompanying derivation
\texttt{DERIVATIVE\_REMAINDER.tex}, based on the original Gamma
beta-limit series, proves
\[
 yL(y)=\alpha_\Gamma y+\frac12+\frac1{8y^2}+E(y),\qquad
 |E(y)|\le\frac1{2y^4}\quad(y>0).
 \tag{LER2}
\]
No asymptotic notation substitutes for this global finite estimate.

\paragraph{Exact moment recursion and its positive comparison data.}
For every integer $p\ge4$, integrate the derivative of
$y^{p+1}\sigma(y)$ on $(0,\infty)$. Its boundary values are zero
because the density is bounded at zero and has the exponential bound
proved in HCT1. The differentiated integrand is integrable by RWB2.
Substitution of LER2 gives the exact identity
\[
 \alpha_\Gamma I_{p+1}=(p+1/2)I_p-\frac18 I_{p-2}-e_p,
 \qquad e_p=\int_0^\infty y^p E(y)\sigma(y)\,dy,
 \qquad |e_p|\le\frac12 I_{p-4}.
 \tag{LER3}
\]
The lower powers on the right are integrable at zero for $p\ge4$.
We retain the sign of $e_p$; the bound on its absolute value is the
only replacement made in an inequality.

RWB11 already gives, with its complete proof from the original density,
\[
 \frac{2v-1}{\pi}\le\frac{I_{v+1}}{I_v}
 \le\frac{2v+3}{\pi}\quad(v\ge0).
 \tag{LER4}
\]
The lower bound is used only when it is positive. For $p\ge10$ define
\[
 U_{p,2}=\frac{\pi^2}{(2p-5)(2p-3)},\qquad
 U_{p,4}=\frac{\pi^4}{(2p-9)(2p-7)(2p-5)(2p-3)},
\]
\[
 b_p=\frac{U_{p,2}/8+U_{p,4}/2}{p+1/2}.
 \tag{LER5}
\]
Telescoping the respective two or four positive lower bounds in LER4
proves $I_{p-2}/I_p\le U_{p,2}$ and
$I_{p-4}/I_p\le U_{p,4}$.
Put
\[
 A_p=\frac{p+1/2}{\alpha_\Gamma},\qquad
 u_p=\frac{I_{p-2}/(8I_p)+e_p/I_p}{p+1/2}.
\]
Then LER3 gives the exact positive ratio
\[
 \frac{I_{p+1}}{I_p}=A_p(1-u_p),\qquad
 |u_p|\le b_p\le\frac7{4p^3}<\frac1{100}\quad(p\ge10).
 \tag{LER6}
\]
For the numerical estimate, $\pi<22/7$ implies $\pi^2<10$ and
$\pi^4<100$. Every factor $2p-9,2p-7,2p-5,2p-3$ is at least $p$
when $p\ge10$. Hence
$b_p\le5/(4p^3)+50/p^5\le7/(4p^3)$, proving every stated constant.

\paragraph{The logarithmic step with a fifth-order remainder.}
For $p\ge12$ define
\[
 V_p=\frac{\alpha_\Gamma^2}{(p-3/2)(p-1/2)},\qquad
 z_p=\frac{V_p}{8(p+1/2)}
     =\frac{\pi^2}{32(p-3/2)(p-1/2)(p+1/2)}.
 \tag{LER7}
\]
There is an exact real $\eta_p$ satisfying
\[
 \boxed{\displaystyle
 \log\frac{I_{p+1}}{I_p}
 =\log\frac{p+1/2}{\alpha_\Gamma}-z_p+\eta_p,
 \qquad |\eta_p|\le\frac{51}{p^5}.}
 \tag{LER8}
\]
To prove it without assuming a moment expansion, multiply the two
identities LER6 at $p-2,p-1$ and invert their positive product:
\[
 \frac{I_{p-2}}{I_p}
 =\frac{V_p}{(1-u_{p-2})(1-u_{p-1})}.
 \tag{LER9}
\]
For $|u|\le b\le1/100$, $|v|\le c\le1/100$, direct subtraction gives
\[
 \left|\frac1{(1-u)(1-v)}-1\right|
 \le\frac{b+c+bc}{(1-b)(1-c)}\le2(b+c).
\]
Indeed $bc\le(b+c)/200$ and $(1-b)(1-c)\ge9801/10000$;
$10050/9801<2$ proves the last bound. LER6 and LER9 thus imply
\[
 \left|\frac{I_{p-2}}{I_p}-V_p\right|
 \le2V_p(b_{p-2}+b_{p-1})\le\frac{7V_p}{(p-2)^3}.
 \tag{LER10}
\]
Consequently the exact $u_p$ satisfies
\[
 |u_p-z_p|\le
 \frac{7V_p/[8(p-2)^3]+U_{p,4}/2}{p+1/2}.
 \tag{LER11}
\]
For $|u|<1$, integration of the derivative of $\log(1-u)+u$ gives
\[
 |\log(1-u)+u|\le\frac{u^2}{2(1-|u|)}.
\]
Combining this with LER6 and LER11 bounds $|\eta_p|$ by
\[
 \frac{7V_p/[8(p-2)^3]+U_{p,4}/2}{p+1/2}
 +\frac{b_p^2}{2(1-b_p)}.
 \tag{LER12}
\]
Here are explicit estimates for the advertised simpler constant.
For $p\ge12$, $p-3/2\ge7p/8$, $p-1/2\ge23p/24$, and $\pi^2<10$
give $V_p<3/p^2$. Also $p-2\ge5p/6$. The first term in LER12 is
therefore bounded by
\[
 \frac{21}{8}\left(\frac65\right)^3p^{-6}+50p^{-5}
 <5p^{-6}+50p^{-5}.
\]
Using LER6, the second term is at most
$[49/(16p^6)](50/99)<2p^{-6}$.
Thus LER12 is less than
$50p^{-5}+7p^{-6}<51p^{-5}$, which proves LER8.

\paragraph{The original scalar endpoint and an exact telescoping correction.}
For integers $a\ge6$ and $l\ge1$ define
\[
 B(x)=\frac1{(x-3/2)(x-1/2)}\quad(x>3/2),\qquad
 K_{a,l}=\frac{\pi^2}{64}\bigl[B(2a)-B(2a+2l)\bigr]>0.
 \tag{LER13}
\]
Summing LER8 at every original integer $p=2a,\ldots,2a+2l-1$ gives
\[
 \boxed{\displaystyle
 \log\frac{\mu_{a+l}}{\mu_a}
 =\sum_{j=0}^{2l-1}\log\frac{2a+j+1/2}{\pi/2}
       -K_{a,l}+\varepsilon_{a,l},\qquad
 |\varepsilon_{a,l}|\le\frac{51l}{16a^5}.}
 \tag{LER14}
\]
In particular no square-pushforward index is substituted for an
all-power moment index: $\mu_{a+l}/\mu_a=I_{2a+2l}/I_{2a}$
is precisely the telescoped ratio. The rational term telescopes
because its exact partial fraction is
\[
 z_p=\frac{\pi^2}{64}\bigl[B(p)-B(p+1)\bigr].
 \tag{LER15}
\]
The sum of the error bounds in LER8 is at most
$51(2l)/(2a)^5=51l/(16a^5)$.
All constants in LER14 are therefore uniform over $a=q,\ldots,q+l$,
as well as over every other integer $a\ge6$.

At $l=1$, this gives the requested individual Christoffel endpoint:
\[
 \boxed{\displaystyle
 \log c_0^{(a)}
 =\log\frac{(4a+1)(4a+3)}{\pi^2}
 -\frac{\pi^2}{64}\bigl[B(2a)-B(2a+2)\bigr]
 +\varepsilon_{a,1},\qquad
 |\varepsilon_{a,1}|\le\frac{51}{16a^5}.}
 \tag{LER16}
\]
The full original low moment remains the product of these coefficients;
the corrections in LER16 add to the single term in LER14, with no
loss from bounding the exact telescoping contribution.

\paragraph{Explicit finite expansion along the original packet.}
Retain $k=4l+1$, $e=1+k(m-1)$ and
$q=e(k+1)^2=e(4l+2)^2$, with $l,m\ge1$. Thus $q\ge36$.
Set
\[
 F_{q,l}=2l\log\frac{4q}{\pi}
 +\frac{l^2}{q}-\frac{l(16l^2-1)}{48q^2}
 +\frac{l^2(8l^2-1)}{48q^3},
\]
\[
 T_{q,l}=\frac{l(768l^4-160l^2+7)}{7680q^4}>0,
 \qquad E_{q,l}=\frac{51l}{16q^5}.
 \tag{LER17}
\]
There is an exact $r_{q,l}$ with $-T_{q,l}\le r_{q,l}\le0$ such that
\[
 \boxed{\displaystyle
 \log\frac{\mu_{q+l}}{\mu_q}
 =F_{q,l}-K_{q,l}+r_{q,l}+\varepsilon_{q,l},\qquad
 |\varepsilon_{q,l}|\le E_{q,l}.}
 \tag{LER18}
\]
Indeed for every $x\ge0$,
\[
 \log(1+x)=x-\frac{x^2}{2}+\frac{x^3}{3}
              -\int_0^x\frac{t^3}{1+t}\,dt,
\]
and the last term belongs to $[-x^4/4,0]$.
Apply this exact formula to $x_j=(j+1/2)/(2q)$, $0\le j<2l$,
in LER14. The required literal power sums, for $M=2l$, are
\[
 \sum_{j=0}^{M-1}(j+1/2)=\frac{M^2}{2},\qquad
 \sum_{j=0}^{M-1}(j+1/2)^2=\frac{M(4M^2-1)}{12},
\]
\[
 \sum_{j=0}^{M-1}(j+1/2)^3=\frac{M^2(2M^2-1)}8,\qquad
 \sum_{j=0}^{M-1}(j+1/2)^4
 =\frac{M(48M^4-40M^2+7)}{240}.
 \tag{LER19}
\]
Each identity follows by expanding the left-hand summand and summing
integer powers; alternatively its right-hand difference at $M+1,M$
equals $(M+1/2)^r$ and its value at $M=0$ is zero, proving it by
induction. Substitution proves every coefficient and the sign in LER18.
The additional rational correction has the elementary bound
\[
 0<K_{q,l}=\sum_{p=2q}^{2q+2l-1}z_p
 \le\frac{3l}{32q^3},
 \tag{LER20}
\]
since the established $V_p<3/p^2$ gives $z_p<3/(8p^3)$.

\paragraph{Exact insertion into the entire Gamma functional.}
Keep the original paired high determinant
$Z_a=H_n^{(a)}H_{n+1}^{(a)}(H_n^{(a+1)})^2$ with $n=q/2$, and
the original product term $P_k$ from LET19. Write
\[
 \mathcal V_k^{\mathrm{high}}=P_k+\log\frac{Z_{q+l}}{Z_q},\qquad
 W_k=\mathcal V_k^{\mathrm{high}}-\log\frac{\mu_{q+l}}{\mu_q}.
 \tag{LER21}
\]
These are the literal original determinant terms, with all four
high blocks retained. LER18 proves the finite receiving interval
\[
 \boxed{\displaystyle
 \mathcal V_k^{\mathrm{high}}-F_{q,l}+K_{q,l}-E_{q,l}
 \ \le W_k\le\
 \mathcal V_k^{\mathrm{high}}-F_{q,l}+K_{q,l}+T_{q,l}+E_{q,l}.}
 \tag{LER22}
\]
In particular the Gamma rational correction is positive in $W_k$,
while the $l^2/q$ term is negative there. The one-sided Taylor error
is also positive after the original low-endpoint minus sign.
This formula updates the scalar contribution only; it leaves the
original high determinants, quotient maps, arithmetic contribution
and their stated bounds intact for their own calculations.

Finally LER17--20 prove, along the actual packet, the more precise
low-endpoint limit
\[
 \log\frac{\mu_{q+l}}{\mu_q}-2l\log\frac{4q}{\pi}
 -\frac{l^2}{q}\longrightarrow0.
 \tag{LER23}
\]
To check this directly, $q\ge(4l+2)^2$ bounds each remaining term by
a constant times respectively $l^{-1},l^{-2},l^{-3},l^{-5},l^{-9}$
as $l\to\infty$, with the displayed exact finite bounds available
in place of these orders. For $m=1$, $l^2/q\to1/16$; for every
fixed $m\ge2$, $l^2/q\to0$ since
$q=[1+(4l+1)(m-1)](4l+2)^2$ grows cubically in $l$.
These conclusions concern the exact original scalar moment ratio;
the complete $W_k$ continues to include both original high ranks.
